\chapter{Caso de estudio}
\label{cap:caso-estudio}

Las vulnerabilidades del capítulo~\ref{cap:explotacion} se han analizado
hasta ahora de forma individual. Sin embargo, en un incidente real, un atacante rara vez se
detiene tras explotar un único fallo: encadena varios, cada uno de
severidad moderada por separado, hasta alcanzar un compromiso total del
sistema. Este patrón de encadenar técnicas individuales hasta lograr un
objetivo (persistencia, movimiento lateral, exfiltración) es precisamente el
que describen catálogos de referencia como MITRE ATT\&CK~\citelib{ATTACK},
aunque este capítulo se limita a un encadenamiento simplificado y no a un
mapeo formal a dichas tácticas y técnicas. Este capítulo reproduce, sobre
\breachlab en modo vulnerable, un
escenario de ataque \emph{end-to-end} que combina cuatro de las seis
vulnerabilidades analizadas.

El escenario parte de un atacante sin ninguna credencial válida y termina
con el control de dos cuentas de usuario distintas, ilustrando cómo la
severidad acumulada de la cadena supera claramente la de cualquiera de sus
eslabones individuales.

\section{Cadena de explotación}

\begin{figure}[htbp]
  \centering
  \begin{tikzpicture}[
      node distance=7mm,
      stepbox/.style={draw, rounded corners, fill=gray!8, minimum width=11.2cm,
        minimum height=9mm, align=left, font=\small, inner sep=3mm},
      arr/.style={-{Latex[length=2mm]}}
    ]
    \node[stepbox] (s1) {\textbf{1. Enumeración} — distingue usuarios válidos por el mensaje de error};
    \node[stepbox, below=of s1] (s2) {\textbf{2. Fuerza bruta} — obtiene la contraseña de \texttt{pablo} (\texttt{qwerty})};
    \node[stepbox, below=of s2] (s3) {\textbf{3. XSS} — roba la cookie de sesión de \texttt{marina}};
    \node[stepbox, below=of s3] (s4) {\textbf{4. IDOR} — lee notas ajenas con la sesión robada};
    \node[stepbox, below=of s4] (s5) {\textbf{5. CSRF} — sobrescribe la contraseña de \texttt{marina} (persistencia)};

    \draw[arr] (s1) -- (s2);
    \draw[arr] (s2) -- (s3);
    \draw[arr] (s3) -- (s4);
    \draw[arr] (s4) -- (s5);
  \end{tikzpicture}
  \caption{Cadena de explotación completa, de la enumeración inicial sin
    credenciales al control persistente de dos cuentas.}
  \label{fig:cadena}
\end{figure}

La figura~\ref{fig:cadena} resume los cinco pasos de la cadena,
detallados a continuación:

\begin{enumerate}
\item \textbf{Enumeración de usuarios (\S\ref{sec:auth}):} probando nombres de usuario
  habituales contra \texttt{POST /api/login}, el atacante distingue por el
  mensaje de error qué cuentas existen (\texttt{pablo}, \texttt{marina},
  \texttt{admin}) sin necesidad de acertar ninguna contraseña.
\item \textbf{Toma de la cuenta de \texttt{pablo} (\S\ref{sec:auth}, fuerza bruta):} al
  no existir límite de intentos, un diccionario básico de contraseñas
  comunes revela que la contraseña de \texttt{pablo} es \texttt{qwerty} en
  el primer intento significativo.
\item \textbf{Robo de la sesión de \texttt{marina} (\S\ref{sec:xss}, XSS):} autenticado
  ya como \texttt{pablo}, el atacante crea una nota con el payload
  \verb|<img src=x onerror="alert(document.cookie)">| y, mediante
  ingeniería social, convence a \texttt{marina} de que abra dicha nota (por
  ejemplo, indicando que contiene información relevante); al no ser la
  cookie \texttt{HttpOnly}, el atacante obtiene el identificador de sesión
  de \texttt{marina}.
\item \textbf{Acceso a información ajena mediante IDOR (\S\ref{sec:idor}):} utilizando la
  sesión de \texttt{marina}, el atacante solicita notas de otros usuarios
  cambiando el identificador en \texttt{/api/notes/<id>}, obteniendo datos
  que ni siquiera \texttt{marina} debería poder ver de otros usuarios.
\item \textbf{Persistencia mediante CSRF (\S\ref{sec:csrf}):} por último, el atacante
  aloja \texttt{csrf\_poc.html} en un sitio externo; si consigue que la
  propia víctima \texttt{marina} (con sesión ya activa tras el paso
  anterior) lo visite, su contraseña se sobrescribe sin su consentimiento,
  garantizando al atacante acceso persistente a la cuenta incluso si la
  sesión robada expira.
\end{enumerate}

\section{Impacto acumulado}

De forma individual, la vulnerabilidad de partida (5.4, fallos de
autenticación) se valoró como severidad media (5.3). Sin embargo, la cadena
completa permite a un atacante externo, sin credenciales iniciales, terminar
con control persistente sobre dos cuentas y acceso de lectura a los datos de
todos los usuarios del sistema, un resultado equiparable en impacto al de la
vulnerabilidad más crítica del capítulo~\ref{cap:explotacion} (la inyección
SQL, 9.1). Esta
discrepancia entre la severidad de cada eslabón por separado y la del ataque
completo es el argumento principal para no evaluar nunca una
vulnerabilidad de forma aislada del resto del sistema.

\section{Efecto de las contramedidas}

Repitiendo la cadena con \texttt{?safe=1} en todos los endpoints, el ataque
se rompe ya en el segundo eslabón: el mensaje deja de permitir enumerar
usuarios y, aunque se conociera un usuario válido, la cuenta se bloquearía
tras 5 intentos. Aun logrando autenticarse por otra vía, el XSS se
neutraliza al renderizarse como texto plano, el acceso a notas ajenas se
rechaza con 403 y la petición CSRF se rechaza sin un token válido. La
cadena ilustra así el valor de la defensa en profundidad: basta con
corregir un solo eslabón para detener el ataque.
